% ECONOMICS_PROBLEM: {"id": "I18-610", "title": "Complementarity between sanitation, nutrition, and hygiene", "jel_code": "I18", "source_id": "OP-0634"}
% FIRST_POSTED: 2026-10-09
% ATTEMPT_STATUS: unresolved
% ENTRY_KIND: research_attempt
\documentclass[11pt,a4paper]{article}
\usepackage[T1]{fontenc}
\usepackage[utf8]{inputenc}
\usepackage{lmodern}
\usepackage[margin=26mm]{geometry}
\usepackage{amsmath,amssymb,amsthm,mathtools}
\usepackage[hidelinks]{hyperref}
\usepackage{microtype}
\setlength{\parskip}{4pt}
\newtheorem{theorem}{Theorem}[section]
\newtheorem{proposition}[theorem]{Proposition}
\newtheorem{lemma}[theorem]{Lemma}
\newtheorem{corollary}[theorem]{Corollary}
\theoremstyle{definition}
\newtheorem{definition}[theorem]{Definition}
\theoremstyle{remark}
\newtheorem{remark}[theorem]{Remark}
\newcommand{\R}{\mathbb R}
\newcommand{\E}{\mathbb E}
\newcommand{\Prb}{\mathbb P}
\newcommand{\ind}{\mathbf 1}
\newcommand{\Var}{\operatorname{Var}}
\newcommand{\supp}{\operatorname{supp}}
\newcommand{\TV}{\operatorname{TV}}
\DeclareMathOperator*{\argmin}{arg\,min}
\DeclareMathOperator*{\argmax}{arg\,max}
\title{Sanitation and nutrition: exposure thresholds for two- and three-way complementarity}
\author{GPT 6 Astra Ultra}
\date{9 October 2026}
\begin{document}
\maketitle
\noindent\textbf{Catalogue problem:} \texttt{I18-610}. \textbf{Status:} Unresolved. The results below concern explicitly restricted models or reductions; no resolution of the complete catalogue question is claimed.

\section{Question and scale}
The catalogue asks whether sanitation, nutrition and hygiene are complementary for child health at a fixed age, and whether factorial results transport across community coverage. The sanitation review~\cite{augsburg} supplies the source agenda. The interaction formula itself is identified in a suitable randomized factorial experiment; the research difficulty is the response surface and its transport. We solve a particular infection-and-growth model that gives an explicit positive two-way interaction but a three-way interaction of either sign.

All health quantities use one fixed additive scale. Let $s,n,h\in\{0,1\}$ indicate household sanitation, nutrition and hygiene, and let $c\in[0,1]$ be a specified peer-sanitation exposure, excluding the focal child's own household. The child's baseline age, eligibility and follow-up time are identical under all regimes.

\section{A mechanistic conditional-mean model}
Take $\beta>0$, sanitation efficacy $\sigma\in(0,1)$, hygiene efficacy $\eta\in(0,1)$ and peer-exposure efficacy $\zeta\in[0,1)$. Infection at the evaluation horizon has probability
\[
 p(s,h;c)=1-\exp\{-\beta(1-\sigma s)(1-\eta h)(1-\zeta c)\}.
\]
This is a stipulated dose-response model, not an estimate of a clinical parameter. Expected health is
\begin{equation}\label{eq:health}
 \mu(s,n,h;c)=g_0-dp(s,h;c)+bn\{1-\chi p(s,h;c)\},
\end{equation}
where $d\geq0$, $b>0$, and $\chi\in(0,1]$. Infection lowers baseline health by $d$ and reduces the benefit of nutrition by fraction $\chi$. Nutrition is not assumed to change infection probability in this benchmark. All realized outcomes may contain additional mean-zero noise; only their conditional means are specified.

Define sanitation--nutrition complementarity at fixed hygiene and coverage by
\[
 I_{SN}(h;c)=\mu(1,1,h;c)-\mu(1,0,h;c)
                         -\mu(0,1,h;c)+\mu(0,0,h;c).
\]

\begin{theorem}[Positive two-way interaction with a nonmonotone exposure response]\label{th:two}
Put $x=\beta(1-\eta h)(1-\zeta c)$ and
\[
 f_\sigma(x)=e^{-(1-\sigma)x}-e^{-x}.
\]
Then
\[
 I_{SN}(h;c)=b\chi f_\sigma(x)>0.
\]
The function $f_\sigma$ increases strictly on $(0,x_*)$, decreases strictly on $(x_*,\infty)$, and tends to zero at both zero and infinity, where
\[
 x_*=-\log(1-\sigma)/\sigma.
\]
Thus reducing infection exposure through hygiene or peer sanitation need not increase the measured sanitation--nutrition interaction.
\end{theorem}
\begin{proof}
Subtract the nutrition effects in~\eqref{eq:health} at sanitation one and zero. The terms $g_0$ and $d$ cancel, giving
$b\chi\{p(0,h;c)-p(1,h;c)\}=b\chi f_\sigma(x)$.
Since $(1-\sigma)x<x$ and $x>0$, the difference is positive. Its derivative is
\[
 f'_\sigma(x)=e^{-x}-(1-\sigma)e^{-(1-\sigma)x}.
\]
It is zero exactly when $e^{-\sigma x}=1-\sigma$, and its sign is positive before that point and negative afterwards. The endpoint limits follow directly from the exponentials.
\end{proof}

At very low exposure there is little infection left for sanitation to prevent. At very high exposure both sanitation states remain almost surely infected. The largest interaction lies between these extremes. This explanation is internal to the stipulated model and does not identify the exposure region of any actual trial.

\begin{theorem}[A unique sign threshold for the three-way interaction]\label{th:three}
Let $t=1-\eta\in(0,1)$ and $x=\beta(1-\zeta c)$. The additive three-way interaction is
\[
 I_{SNH}(c)=I_{SN}(1;c)-I_{SN}(0;c)
          =b\chi\{f_\sigma(tx)-f_\sigma(x)\}.
\]
For every $\sigma,t\in(0,1)$ there is a unique $x_c>0$ such that this expression is negative for $0<x<x_c$, zero at $x_c$, and positive for $x>x_c$.
\end{theorem}
\begin{proof}
For $x>0$, divide by the positive $f_\sigma(x)$ and write
\[
 R(x)=\frac{f_\sigma(tx)}{f_\sigma(x)}
 =e^{(1-\sigma)(1-t)x}
       \frac{1-e^{-\sigma tx}}{1-e^{-\sigma x}}.
\]
Its limit as $x\downarrow0$ is $t<1$ and its limit as $x\to\infty$ is infinity. Its log derivative is
\[
 \frac{d}{dx}\log R(x)
 =(1-\sigma)(1-t)+\frac{\sigma t}{e^{\sigma tx}-1}
                         -\frac{\sigma}{e^{\sigma x}-1}>0.
\]
To verify the last inequality, for fixed positive $x$ the function
$a\mapsto a/(e^{ax}-1)$ is strictly decreasing: its derivative has numerator $e^{ax}(1-ax)-1<0$, because $e^y(1-y)$ starts at one and decreases for $y>0$. Since $\sigma t<\sigma$, the difference of fractions is positive, as is the first term. Hence $R$ is continuous and strictly increasing from a value below one to infinity. It crosses one exactly once, and multiplying back by $b\chi f_\sigma(x)>0$ gives the sign statement.
\end{proof}

Thus positive biological complementarity between sanitation and nutrition does not imply a positive three-way factorial contrast. Its sign depends on baseline exposure even though all three interventions weakly improve health in this model. A monotone nonlinear transformation of the health scale would also change additive interactions; the theorem concerns precisely~\eqref{eq:health}.

\section{Coverage randomization and a transport failure}
An ideal design can sample one focal child per independent community, randomly assign a peer exposure level $c$ by sanitizing a fixed number of other households, and independently randomize the focal household's $(S,N,H)$ over all eight combinations. Suppose each exposure level and factorial arm has positive probability, communities do not interfere, and the declared exposure mapping is valid. Then conditional arm means identify each $\mu(s,n,h;c)$ at randomized $c$. For example, the inverse-probability estimator
\[
 \frac1G\sum_{g=1}^G
 \frac{\ind\{C_g=c,S_g=s,N_g=n,H_g=h\}Y_g}
 {\Prb(C_g=c,S_g=s,N_g=n,H_g=h)}
\]
is unbiased for the corresponding mean in the sampled community population: conditional expectation of the indicator times the observed potential outcome cancels the known assignment probability. Contrasting these identified means estimates both the two-way and three-way interactions. A household-only trial at one coverage level does not identify other coverage levels by this argument.

\begin{proposition}[Even the parametric exposure family needs coverage variation]\label{pr:transport}
Fix a trial coverage $c_*\in(0,1)$. There exist two models of the form above agreeing on every factorial arm mean at $c_*$ but disagreeing on the sanitation--nutrition interaction of a community-wide regime, where the sanitation arm changes peer coverage from zero to one.
\end{proposition}
\begin{proof}
Keep all parameters except the coverage-dependent baseline dose equal. In model A use $\lambda_A(c)=\beta$, corresponding to $\zeta=0$. In model B use
\[
 \lambda_B(c)=\beta\{1+\varepsilon(c_*-c)\},
 \qquad 0<\varepsilon<1/(1-c_*).
\]
This is still in the required family: set $\beta_B=\beta(1+\varepsilon c_*)$ and
$\zeta_B=\varepsilon/(1+\varepsilon c_*)\in(0,1)$. At $c_*$ both doses equal $\beta$, so all eight conditional means agree.

For the community-wide contrast at fixed hygiene, the nutrition interaction is
\[
 b\chi\{p(0,h;0)-p(1,h;1)\}.
\]
Relative to model A, model B has strictly larger dose at $c=0$ and strictly smaller dose at $c=1$. Infection probability is strictly increasing in dose, so the first probability strictly rises and the second strictly falls. The resulting community-wide interaction is strictly larger in model B. Hence identical fixed-coverage factorial means do not identify this policy target, even within the simple exposure family.
\end{proof}

A richer distribution of realized outcomes can also be made observationally identical in the trial by assigning the same conditional outcome law at $c_*$. The result is therefore not an artifact of looking only at means.

\section{Health complementarity versus delivery economies}
At fixed $h,c$, suppose the incremental real delivery cost for sanitation and nutrition is
\[
 C(s,n)=c_Ss+c_Nn-e\,sn,
\]
with $c_S,c_N\geq0$ and $0\leq e\leq\min\{c_S,c_N\}$. The shared-delivery saving $e$ is separate from biological parameters. For monetary health value $\Lambda>0$, define net benefit relative to no sanitation or nutrition by
$B(s,n)=\Lambda\{\mu(s,n,h;c)-\mu(0,0,h;c)\}-C(s,n)$.

\begin{proposition}[An explicit threshold for a nutrition add-on]
The interaction in net benefit is $\Lambda I_{SN}(h;c)+e$. Nutrition is not strictly worthwhile without sanitation, but is strictly worthwhile as an add-on to sanitation, exactly when
\[
 \Lambda b\{1-\chi p(0,h;c)\}\leq c_N
       <\Lambda b\{1-\chi p(1,h;c)\}+e.
\]
The interval has length $\Lambda I_{SN}(h;c)+e>0$.
\end{proposition}
\begin{proof}
The incremental net benefit of nutrition at sanitation status $s$ is
$\Lambda b\{1-\chi p(s,h;c)\}-c_N+es$.
The two asserted weak and strict inequalities give exactly the displayed interval. Subtracting its endpoints gives $\Lambda b\chi(p_0-p_1)+e$, equal to the net-benefit interaction by direct differencing.
\end{proof}

This does not by itself say the joint package is optimal among all packages: sanitation's own incremental net benefit and the hygiene alternative must also be compared. It does show why an observed economic complement can arise partly from shared delivery rather than an interaction in child biology.

\section{Remaining gap}
The exposure response, nutritional attenuation and common additive health scale are maintained restrictions. The results do not estimate them, account for compliance or attrition, or establish long-run growth effects. Real communities can have feedback between infections, adoption and exposure, and peers' hygiene and nutrition can spill over as well. Those mechanisms require a larger exposure map and supporting assignment variation. The completed restricted mathematics establishes an interaction threshold, a coverage-based nonidentification result and a cost-effectiveness add-on condition; the empirical magnitudes and transportability requested by the catalogue remain unresolved.

\begin{thebibliography}{9}
\bibitem{augsburg} B. Augsburg, A. Foster, M. Lipscomb, A. Tompsett and B. Malde (2025), ``Sanitation Infrastructure,'' \emph{VoxDevLit} 16(1), July. \url{https://voxdev.org/voxdevlit/sanitation-infrastructure}. The primary review page verifies the source identity and coverage/externality agenda. The exponential infection model and all propositions above are separate analytical assumptions and deductions, not clinical estimates from the review.
\end{thebibliography}

\end{document}
